\scanpage{35}
\mspage{5}
allusion, on devrait pouvoir exprimer \inserted{alors}
le \uline{type d'homotopie} \uline{de} $\deleted{\unsure{Z}}X$
(et non seulement son $\pi_1$) en termes du
diagramme \inserted{de groupoïdes} (8), ou ce qui revient au
$=$, du diagramme de groupes (10).
En tous cas, il est clair (indépendam-
ment de toute hypothèse de nullité de
certains $\pi_i$, ou de locale trivialité normale)
comment reconstruire, en termes du
diagramme 8, la \inserted{\unsure{ss}} catégorie \inserted{$\unsure{\mathcal F|\mathfrak F}$} de la cat. des
faisceaux sur $X$, \deleted{dont} formée des $F$ tels
que l'on ait
\[
(16)\qquad F\mid X^*\quad\text{et}\quad F\mid Y\quad\text{localement triviaux},
\]
Cette catégorie \inserted{$\unsure{\mathcal F|\mathfrak F}$} est équivalente en effet
à celle des systèmes
\[
(17)\qquad (E_{X^*},E_{Y,X};\varphi)
\]
\deleted{\unsure{formés}}\inserted{où $E_{X^*}$ est un}
\deleted{d'un} système local sur $\Pi_1 X^*$\inserted{($\unsure{\leftrightarrow|\rightarrow}$ un rev. \inserted{étale} de $X^*$)}, $E_{Y,X}$ un
syst. local sur $\Pi_1 V_{Y,X}$\inserted{$\xleftarrow{\approx}\Pi_1Y$} ($\unsure{\leftrightarrow|\rightarrow}$ un rev. étale
de $Y$) et $\varphi$ un hom. de syst. locaux sur $\Pi_1V^*_{Y,X}$
\mspage{34}
